\documentclass{article}
\usepackage{amssymb, amsmath}

\title{Continued Fractions}
\author{Nick Wendt}
\date{December 18, 2025}

\begin{document}
\maketitle
\section*{Definition}
A simple continued fraction is an expresion of the form
$$
a_0+\frac{1}{a_1+\frac{1}{a_2+\frac{1}{a_3+...}}}
$$
Where $a_0 \in \mathbb{Z}$ and $a_1, ...,a_n \in \mathbb{N}$
\\
We write a continued fraction compactly as $[a_0;a_1,a_2,a_3,...]$
\section*{Simple Example}
Write $\frac{31}{12}$ as a continued fraction
Step 1: Find the floor of $\frac{31}{12}$
$$
\left\lfloor{\frac{31}{12}}\right\rfloor = 2
$$
This is your $a_0$\\
\\
Step 2: Subtract your original number by the floor calculated in step 1.
$$
\frac{31}{12} - 2 = \frac{31}{12} - \frac{24}{12} = \frac{7}{12}
$$
This will always give you a number between 0 and 1.\\
\\
Step 3: Flip the fraction 
$$
\frac{7}{12} \rightarrow \frac{12}{7}
$$
Step 4: Repeat steps 1-3 until you get a 1 as the numerator of the fraction in step 2.\\
\\
Finishing the example:
\begin{gather*}
Floor:\ \left\lfloor{\frac{12}{7}} \right\rfloor = 1 = a_1\\
Subtract:\ \frac{12}{7} - 1 = \frac{5}{7}\\
Flip:\ \frac{5}{7} \rightarrow \frac{7}{5}\\
Repeat\\
Floor:\  \left\lfloor{\frac{7}{5}}\right\rfloor = 1 = a_2\\
Subtract:\ \frac{7}{5} - 1 = \frac{2}{5}\\
Flip:\ \frac{2}{5} \rightarrow \frac{5}{2}\\
Floor:\ \left\lfloor{\frac{5}{2}}\right\rfloor = 2 = a_3\\
Subtract:\ \frac{5}{2} - 2 = \frac{1}{2}\\
Stop
\end{gather*}
We stop since 1 is in the numerator of the fraction.\\
\\
Now using our $a_i's$, we can write our continued fraction:
$$
\frac{31}{12} = 2 + \frac{1}{1+\frac{1}{1+\frac{1}{2+\frac{1}{2}}}}
$$

\section*{Irrational Example}

\end{document}